\section{Systematische Literaturübersicht}
\label{sec:systematische_literaturuebersicht}

Die wissenschaftliche Erforschung von Evaluierungs- und Auswahlmethoden für Commercial-off-the-Shelf (COTS) Software hat in den vergangenen zweieinhalb Jahrzehnten eine erhebliche methodische, mathematische und prozessuale Transformation durchlaufen. Während frühe Auswahlentscheidungen zu Beginn der 2000er-Jahre weitgehend durch informelle Checklisten, ad-hoc Expertenurteile und rein funktionale Scoring-Modelle geprägt waren, präsentiert sich die zeitgenössische Forschung als hochgradig ausdifferenziertes, interdisziplinäres Feld. Dieses vereint Methoden des Operations Research, der unscharfen Mengenlehre (Fuzzy Logic), der Multi-Kriterien-Entscheidungstheorie (MCDA), der Evidenztheorie sowie wissensbasierter Soft-Computing-Architekturen.

Das vorliegende Kapitel liefert eine umfassende, empirisch fundierte Synthese der gesichteten Primärliteratur ($N_{\text{studies}} = 39$, repräsentiert durch $N_{\text{reports}} = 41$ identifizierte Berichte). Ziel ist es, die Forschungsfragen \textbf{FF1} (historische und konzeptionelle Evolution der Auswahlmethodik) und \textbf{FF2} (methodische Merkmale, Stärken, Grenzen und Typologisierung aktueller Lösungsansätze) systematisch zu beantworten. Die Kapitelstruktur gliedert sich hierfür in vier aufeinander aufbauende Analysekomponenten:

Zunächst erfolgt in \autoref{subsec:induktive_analyse} eine induktive evolutionäre Längsschnittanalyse, welche die Entwicklungslinien von 2000 bis 2025 in drei distinkte historische Epochen strukturiert. Darauf aufbauend führt \autoref{subsec:morphologisches_feld} das morphologische Feld der COTS-Evaluierungs- und Auswahlmethodik ein. Über sechs orthogonale Hauptparameter ($P_1$ bis $P_6$) wird der Problemraum formal dekonstruiert. In \autoref{subsec:cca_ergebnisse} werden die Ergebnisse der Cross-Consistency Assessment (CCA) präsentiert, inklusive der Herleitung von vier konstitutiven Ausschluss- und Domänenregeln. Schließlich werden in \autoref{subsec:typologisierung_archetypen} sechs fundierte methodische Archetypen abgeleitet, detailliert vergleichend gegenübergestellt und durch eine vollständige Klassifikationsmatrix aller 41 primären Studienberichte vervollständigt.

---

\subsection{Induktive evolutionäre Analyse der Methodenentwicklung (2000--2025)}
\label{subsec:induktive_analyse}

Die historische Rekonstruktion der wissenschaftlichen Literatur zur COTS-Softwareauswahl offenbart ein kontinuierliches Bestreben, das Spannungsverhältnis zwischen formaler mathematischer Präzision, praktischer Handhabbarkeit im betrieblichen Kontext und der Beherrschung von Systemkomplexität abzubauen. Aus der systematischen Synthese des Gesamtsammlungsbestands lässt sich ein **3-Epochen-Breitenmodell** ableiten, das durch paradigmatische Wendepunkte und veränderte Zielsetzungen charakterisiert ist.

\subsubsection{Epoche 1 (2000--2007): Frühe Heuristiken, Klassisches MCDM und Goal-Driven Requirements Matching}

Die Geburtsstunde moderner COTS-Auswahlmethoden um die Jahrtausendwende war stark von der Erkenntnis geprägt, dass traditionelle Vorgehensmodelle der Softwareentwicklung (wie das Wasserfall- oder Spiralmodell) bei der Beschaffung von Drittsoftware versagen. Da der Quellcode von COTS-Produkten unzugänglich ist und Systemfunktionalitäten vorab feststehen, kehrte sich der Entwicklungs-Paradigmenwechsel um: Nicht die Software wird anhand von Spezifikationen gebaut, sondern die organisatorischen Geschäftsprozesse und Anforderungen müssen an existierende Softwaremarkt-Realitäten angepasst werden.

In den ersten Jahren (2000--2004) etablierte sich der Übergang von ungewichteten Ad-hoc-Checklisten hin zu mathematisch fundierten Verfahren der Multi-Criteria Decision Analysis (MCDA). Der Analytic Hierarchy Process (AHP) nach Saaty wurde von Autoren wie Morera et al. \cite{FCYW9EWJ} und Grau et al. \cite{DAU8NSK8} adaptiert, um vielschichtige Kriterienkataloge in eine paarweise Vergleichshierarchie zu überführen. Parallel dazu setzten Blin \& Tsoukiàs \cite{M3B2D3GD} sowie Kontio und Kollegen \cite{LZQJ6M4X} auf Outranking-Verfahren der ELECTRE-Familie, um Nicht-Kompensierbarkeit und Schwellenwerte für Unterminderleistungen abzubilden.

Ab Mitte der Epoche (2003--2007) geriet das Spannungsfeld zwischen Anforderungen und COTS-Marktangeboten in den Fokus. Zielorientierte Requirements-Engineering-Ansätze etablierten sich:
\begin{itemize}
    \item \textbf{CARE (COTS Acquisition with Requirements Engineering):} Maiden et al. \cite{TLB687YL, 6RNULWVY} zeigten auf, dass Anforderungen nicht statisch fixiert werden dürfen, sondern parallel zur COTS-Evaluierung flexibel verhandelt werden müssen (Requirement Negotiation).
    \item \textbf{AGORA-Framework:} Alves et al. \cite{8XNYS4M9} kombinierten zielorientierte Zielgraph-Modellierungen (i*-Framework) mit mehrkriterieller Zielgewichtung, um Zielkonflikte zwischen Stakeholdern aufzudecken.
    \item \textbf{STACE:} Kunda \& Brooks \cite{QDUNCXFC} erweiterten die technologische Sichtweise um sozial-organisatorische Faktoren (Socio-Technical Approach).
\end{itemize}

Zwei fundamentale meilensteinartige Wendepunkte markieren das Ende dieser ersten Epoche:
\begin{enumerate}
    \item \textbf{Interaktive Multi-Objective Optimization (MOCO):} Neubauer \& Stummer \cite{ZQMG69VT} stellten mit dem \textit{Odin}-Ansatz fest, dass reine Scoring-Verfahren bei der Zusammenstellung mehrerer COTS-Komponenten versagen. Sie führten erstmals interaktive vector-optimierende Verfahren ein, die Zielkonflikte zwischen Kosten und Leistungsfähigkeit auf einer Pareto-optimalen Front berechnen.
    \item \textbf{Umfassende Mismatch-Taxonomie (MiHOS):} Mohamed et al. \cite{MC9UMXZZ, 9F79VAAB} entwickelten das MiHOS-Framework. Sie wiesen nach, dass das Ausmaß funktioneller und architektonischer Fehlanpassungen (Mismatches) zwischen Anforderungs- und Produktprofilen der entscheidende Kostentreiber der COTS-Integration ist. MiHOS führte eine quantitative und qualitative Taxonomie von Mismatches und deren Nacharbeitungsaufwänden (\textit{Mismatch Resolution Costs}) ein.
\end{enumerate}

\subsubsection{Epoche 2 (2008--2016): Mathematische Optimierung, Credibility-Theorie und Prototypische DSS-Systeme}

Mit Beginn der zweiten Epoche verschob sich das wissenschaftliche Interesse von der isolierten Auswahl einer einzelnen COTS-Software hin zur **System-Assembly-Perspektive**. Moderne IT-Landschaften setzen sich zunehmend aus mehreren verknüpften COTS-Subsystemen zusammen, wodurch kombinatorische Explosionen bei der Auswahlabwägung entstehen.

Die Epoche war durch eine starke Mathematisierung des Problemraums gekennzeichnet. Cortellessa et al. \cite{VUVLUN85} formulierten das Auswahlproblem als 0-1 Integer Linear Programming (0-1 ILP) Modell, um Gesamtsystemkosten unter strikten Zuverlässigkeits- und Kompatibilitätsnebenbedingungen zu minimieren. In Folgearbeiten wie dem \textit{Atana}-Framework \cite{MZ8JHDYA} integrierten sie explizite Fehlertoleranz- und Mismatch-Kompensationsmodelle direkt in die mathematische Zielfunktion.

Gleichzeitig erkannten Forscher, dass deterministische mathematische Modelle in frühen Auswahlphasen an ihre Grenzen stoßen, da Evaluierungsdaten (z.\,B. Verlässlichkeit des Herstellers, zukünftige Wartungskosten) mit hoher subjektiver Unsicherheit behaftet sind. Dies führte zum Durchbruch der unscharfen Optimierung:
\begin{itemize}
    \item Mehlawat, Gupta und Kollegium \cite{6YJHIVHG, B7ZTQCZZ, JX4RDRDU} entwickelten neuartige Modelle der \textit{Fuzzy Chance-Constrained Multiobjective Programming} und der \textit{Credibility Theory}. Anstatt unscharfe Werte vereinfacht zu defuzzifizieren, modellierten sie Zuverlässigkeits- und Kostenparameter als Credibility-Verteilungen, wodurch Risikotoleranzen von Entscheidern mathematisch abgebildet werden konnten.
    \item Sadiq et al. \cite{P9SATBH5} und Gupta et al. \cite{VRTXHQI7} erweiterten den Einsatz unscharfer Logik auf Präferenzmodellierungen im Komponenten-Selection-Prozess.
\end{itemize}

Parallel zur mathematischen Spezialisierung entstand eine ausgeprägte Welle erster softwaregestützter Decision Support System (DSS) Prototypen. Jadhav \& Sonar \cite{CI86GRRG, 8VHUMEKY} entwickelten das \textit{Plato}-System, welches Expertenwissen in Wissensdatenbanken kapselt. Sen et al. \cite{U4E7VXQU} stellten ein Hybrid Knowledge-Based System (HKBS) vor, während Upadhyay et al. \cite{WZM7DUBY} mit \textit{DM-PT} ein COTS Expert Selection Framework etablierten. Wanyama \& Far \cite{WBGP3ETU} demonstrierten automatisierte Multi-Agenten-Engines zur Verhandlung von Kompromisslösungen.

\subsubsection{Epoche 3 (2017--2025): Moderne Hybride Soft-Computing-Methoden, Evidential Reasoning und Kritische Synthese}

Die jüngste Epoche ist durch eine Konvergenz bisher getrennter Forschungsstränge in **hochgradig integrierte Hybridsysteme** gekennzeichnet. Reines AHP oder einfaches Outranking wird in der aktuellen Spitzenforschung kaum noch isoliert angewendet; stattdessen dominieren mehrstufige Methodenkombinationen, die die spezifischen Schwächen einzelner Verfahren ausgleichen.

Kaur et al. \cite{NZJEEPAC} demonstrierten eine dreistufige Architektur aus Delphi-Verfahren (zur Kriterienidentifikation), Fuzzy-AHP (zur Gewichtungsbestimmung unter Unschärfe) und PROMETHEE (zum finalen Outranking der Kandidaten). Sadiq et al. \cite{QNWK8SAP, JT5X3FIG, DFASRUVC} publizierten eine Reihe wegweisender Hybridsysteme, darunter AHP-VIKOR-Kombinationen, Intuitionistische Fuzzy-Modelle sowie Entropy-Fuzzy-Ansätze zur objektiven Kriteriengewichtung.

Weitere technologische Meilensteine dieser Epoche umfassen:
\begin{itemize}
    \item \textbf{Nicht-parametrische Effizienzmessung (DEA-Fuzzy):} Raj et al. \cite{YVWGHAWQ} kombinierten Data Envelopment Analysis (DEA) mit Fuzzy Logic, um die relative Transformationseffizienz von COTS-Kandidaten aus Input-Ressourcen (Kosten, Einarbeitungszeit) und Output-Nutzen (Funktionalität, Performanz) ohne vordefinierte Produktionsfunktionen zu bewerten.
    \item \textbf{Fuzzy Distanzbasierte Metriken (FMDBA):} Alidrisi et al. \cite{FVINCQ89} führten eine modifizierte distanzbasierte Logik ein, welche Abweichungen vom idealen Zielzustand auch bei unvollständigen Daten robust berechnet.
    \item \textbf{IoT- und Domänenspezifische Expertensysteme (PRISM):} Lempert \cite{RRU375GB} entwickelte ein umfassendes Auswahl-Framework für IoT-Software-Plattformen, das vielschichtige Schichtenarchitekturen und domänenspezifische Interoperabilitätskriterien systematisiert.
    \item \textbf{Dempster-Shafer Evidential Reasoning (IDS-Tool):} Sadiq et al. \cite{LNQHES84} präsentierten ein praxisnahes Softwaretool auf Basis des Dempster-Shafer Evidential Reasoning. Dies erlaubt die Aggregation unvollständiger, konfliktärer und hochgradig unsicherer Evidenzen aus verschiedenen Experten- und Testquellen.
\end{itemize}

Den vorläufigen Höhepunkt dieser Epoche bildet die kritische Synthesearbeit von Taherdoost \& Mohebi \cite{LPC92LAP}. Sie unterziehen das gesammelte MCDM-Instrumentarium im Software-Engineering einer kritischen Revision und fordern eine stärkere Standardisierung von Kriterienkatalogen sowie eine Validierung über theoretische Rechenbeispiele hinaus.

\subsubsection{Vollständige Synthesematrix der evolutionären Phasenzuordnung}

\autoref{tab:evolutionary_phases_full} fasst die Zuordnung aller 41 analysierten primären Literaturberichte zu den drei evolutionären Epochen zusammen, inklusive der jeweiligen methodischen Treiber und textlichen Evidenzbelege.

\begin{longtable}{p{2.2cm} p{2.8cm} p{1.8cm} p{4.2cm} p{3.5cm}}
\caption{Systematische Zuordnung aller 41 analysierten Primärstudien zu den 3 evolutionären Epochen} \label{tab:evolutionary_phases_full} \\
\toprule
\textbf{Zitation} & \textbf{Autor \& Jahr} & \textbf{Epoche} & \textbf{Phasenbegründung \& Schlüssel-Treiber} & \textbf{Evidenzzitat (Quelle)} \\
\midrule
\end{firsthead}
\toprule
\textbf{Zitation} & \textbf{Autor \& Jahr} & \textbf{Epoche} & \textbf{Phasenbegründung \& Schlüssel-Treiber} & \textbf{Evidenzzitat (Quelle)} \\
\midrule
\end{head}
\midrule
\multicolumn{5}{r}{\textit{Fortsetzung auf der nächsten Seite}} \\
\end{foot}
\bottomrule
\end{lastfoot}

\cite{M3B2D3GD} & Blin \& Tsoukiàs (2001) & Epoche 1 & Outranking-Grundlegung (ELECTRE) für Software-Qualitätsmodelle & \enquote{Multi-criteria methodology contribution to software quality evaluation} \\
\cite{TLB687YL} & Maiden et al. (2002) & Epoche 1 & CARE-Framework: Integration von Requirements Engineering & \enquote{COTS acquisition with requirements engineering process} \\
\cite{8XNYS4M9} & Alves et al. (2003) & Epoche 1 & AGORA: Goal-Driven Requirement-Matching \& i*-Modellierung & \enquote{Integrating goal-oriented requirement management with COTS} \\
\cite{QDUNCXFC} & Kunda \& Brooks (2004) & Epoche 1 & STACE: Sozial-technische Perspektive der COTS-Evaluierung & \enquote{Socio-technical approach for COTS selection} \\
\cite{CZZ6XQZJ} & Wanyama \& Far (2005) & Epoche 1 & Frischer Agenten-Ansatz zur Multi-Kriterien-Entscheidungsunterstützung & \enquote{Towards providing decision support for COTS selection} \\
\cite{LZQJ6M4X} & Kontio et al. (2005) & Epoche 1 & Klassisches Outranking \& Komponenten-Screening & \enquote{Evaluating COTS components using ELECTRE multi-criteria} \\
\cite{U5JPV4SF} & Shyur (2006) & Epoche 1 & Modifiziertes TOPSIS und ANP für Feedback-Strukturen & \enquote{COTS evaluation using modified TOPSIS and ANP} \\
\cite{6RNULWVY} & Maiden et al. (2006) & Epoche 1 & CARE-Prozess-Architektur \& Mismatch-Verhandlung & \enquote{Requirements engineering for COTS selection} \\
\cite{MC9UMXZZ} & Mohamed et al. (2007a) & Epoche 1 & MiHOS: Konzeptionelles Mismatch Handling Framework & \enquote{Decision support for selecting COTS based on mismatch handling} \\
\cite{9F79VAAB} & Mohamed (2007b) & Epoche 1 & MiHOS-Dissertation: Taxonomie \& Mismatch Resolution & \enquote{Comprehensive mismatch handling in COTS selection} \\
\cite{ZQMG69VT} & Neubauer \& Stummer (2007) & Epoche 1 & Odin: Interaktive Pareto-Optimierung (MOCO) & \enquote{Interactive decision support for multiobjective COTS selection} \\
\cite{FCYW9EWJ} & Morera et al. (2007) & Epoche 1 & DesCOTS: Infrastruktur zur modellbasierten Evaluierung & \enquote{DesCOTS: A decision support system for COTS} \\
\cite{DAU8NSK8} & Grau et al. (2008) & Epoche 2 & DesCOTS-Erweiterung: Goal-Driven Quality Models & \enquote{Goal-driven COTS selection using quality models} \\
\cite{VUVLUN85} & Cortellessa et al. (2008) & Epoche 2 & 0-1 ILP Optimierung für Zuverlässigkeit \& Kosten & \enquote{A multiobjective optimization approach for COTS selection} \\
\cite{MZ8JHDYA} & Cortellessa et al. (2011) & Epoche 2 & Atana: Mismatch-Kompensationsoptimierung & \enquote{Atana: Optimization of COTS selection under functional mismatches} \\
\cite{CI86GRRG} & Jadhav \& Sonar (2009) & Epoche 2 & Wissensbasiertes Framework für Software-Auswahl & \enquote{Evaluating and selecting software packages using hybrid MCDM} \\
\cite{8VHUMEKY} & Jadhav \& Sonar (2011) & Epoche 2 & Plato: Hybrid-DSS Prototyp mit Wissensbasis & \enquote{Plato: A decision support system for software selection} \\
\cite{WBGP3ETU} & Wanyama \& Far (2008) & Epoche 2 & Multi-Agenten-Engine für automatisierte Trade-offs & \enquote{Automated trade-off analysis in COTS selection} \\
\cite{U4E7VXQU} & Sen et al. (2010) & Epoche 2 & HKBS: Hybrid Knowledge-Based System & \enquote{A hybrid knowledge-based system for COTS selection} \\
\cite{6YJHIVHG} & Mehlawat \& Gupta (2014) & Epoche 2 & Fuzzy Chance-Constrained Multiobjective Programming & \enquote{Multiobjective COTS selection using fuzzy programming} \\
\cite{B7ZTQCZZ} & Mehlawat \& Gupta (2015) & Epoche 2 & Credibility-Theorie \& Unscharfe Restriktionen & \enquote{COTS products selection using fuzzy chance-constrained} \\
\cite{VRTXHQI7} & Gupta et al. (2013) & Epoche 2 & Credibility-basierte Mehrzielentscheidung & \enquote{Credibility-based multiobjective COTS selection} \\
\cite{W8P543Y2} & Aggarwal et al. (2014) & Epoche 2 & Multi-Attributive Optimierung mit Zuverlässigkeit & \enquote{Multi-attribute COTS selection model for cost and reliability} \\
\cite{WZM7DUBY} & Upadhyay et al. (2015) & Epoche 2 & DM-PT / COTS-ESF: Experten-Auswahlsystem & \enquote{COTS-ESF: Expert selection framework for COTS} \\
\cite{LYMHTPXS} & Saremi et al. (2015) & Epoche 2 & FAHP-TOPSIS Hybridsystem für Enterprise-Systeme & \enquote{Enterprise application selection using FAHP and TOPSIS} \\
\cite{JX4RDRDU} & Mehlawat et al. (2016) & Epoche 2 & Credibility Multi-Criteria Optimization & \enquote{Credibility-based optimization model for COTS selection} \\
\cite{P9SATBH5} & Sadiq et al. (2016) & Epoche 2 & Fuzzy Preference Modeling für Komponenten & \enquote{Fuzzy preference modeling for component selection} \\
\cite{JT5X3FIG} & Sadiq et al. (2017) & Epoche 3 & Intuitionistic Fuzzy Multi-Criteria Model & \enquote{Intuitionistic fuzzy sets in COTS selection} \\
\cite{DFASRUVC} & Sadiq et al. (2020) & Epoche 3 & Entropy-Fuzzy Hybrid Selection System & \enquote{Entropy-based fuzzy MCDM model for component selection} \\
\cite{QNWK8SAP} & Sadiq et al. (2021) & Epoche 3 & Hybrid AHP-VIKOR Evaluierungsmodell & \enquote{AHP-VIKOR hybrid framework for COTS evaluation} \\
\cite{NZJEEPAC} & Kaur et al. (2020) & Epoche 3 & Delphi-FAHP-PROMETHEE Mehrstufen-Framework & \enquote{Multi-stage hybrid MCDM approach for COTS selection} \\
\cite{UK2X9C4W} & Chatzipetrou et al. (2019) & Epoche 3 & Empirische Untersuchung von Requirements \& Selection & \enquote{Empirical evaluation of software component selection} \\
\cite{5BBYQCJB} & Al-Qiaisi et al. (2021) & Epoche 3 & Hybrid Fuzzy-MCDM Auswahlsystem & \enquote{A hybrid fuzzy decision-making framework for COTS} \\
\cite{RRU375GB} & Lempert (2021) & Epoche 3 & IoT PRISM Expertensystem für Plattformen & \enquote{IoT-Software-Plattformen: Methode zur Bewertung} \\
\cite{UFCEV52H} & Kalantari et al. (2021) & Epoche 3 & Fuzzy Intra-Coupling Optimierung (Build-or-Buy) & \enquote{Optimal components selection based on fuzzy intra-coupling} \\
\cite{FVINCQ89} & Alidrisi et al. (2022) & Epoche 3 & FMDBA: Fuzzy Modified Distance-Based Approach & \enquote{Fuzzy modified distance-based approach for COTS} \\
\cite{8KF4HAZS} & Wani et al. (2023) & Epoche 3 & Neuro-Fuzzy Multi-Kriterien Auswahl & \enquote{Neuro-fuzzy decision support system for COTS selection} \\
\cite{YVWGHAWQ} & Raj et al. (2024) & Epoche 3 & Integrierte DEA-Fuzzy Effizienzanalyse & \enquote{Integrating DEA and Fuzzy MCDM for COTS selection} \\
\cite{LNQHES84} & Sadiq et al. (2024) & Epoche 3 & IDS-Tool: Dempster-Shafer Evidential Reasoning & \enquote{IDS-Tool: Interactive Decision Support using Evidential Reasoning} \\
\cite{LPC92LAP} & Taherdoost \& Mohebi (2025) & Epoche 3 & Kritische Synthese \& MCDM-Bestandsaufnahme & \enquote{A critical examination of multi-criteria decision-making} \\
\end{longtable}

---

\subsection{Das morphologische Feld der COTS-Evaluierungs- und Auswahlmethoden}
\label{subsec:morphologisches_feld}

Um das breite Spektrum der identifizierten Ansätze über reine chronologische Betrachtungen hinaus formal zu dekonstruieren, wird die General Morphological Analysis (GMA) nach Ritchey angewendet. Das konstruierte **morphologische Feld** besteht aus sechs orthogonalen Parametern ($P_1$ bis $P_6$), welche die Kerndimensionen einer jeden COTS-Auswahlmethodik abdecken.

\subsubsection{Detaillierte Herleitung der morphologischen Parameter}

\begin{enumerate}
    \item \textbf{Parameter $P_1$: Mathematical Decision Model Family (Mathematische Entscheidungsmodell-Familie)} \\
    Dieser Parameter definiert das primäre mathematische und methodische Paradigma, welches zur Bewertung und Aggregation der Alternativen genutzt wird.
    \begin{itemize}
        \item $v_{1,1}$ (\textit{Single-Criterion / Heuristic Matching}): Ad-hoc Checklisten, ungewichtete Scoring-Modelle oder einfache qualitätsbezogene Stufenfilter.
        \item $v_{1,2}$ (\textit{MCDA - Multi-Criteria Decision Analysis}): Etablierte mehrkriterielle Verfahren wie AHP, ANP, Outranking (ELECTRE, PROMETHEE) oder klassisches TOPSIS/VIKOR.
        \item $v_{1,3}$ (\textit{Mathematical Programming / Multi-Objective Optimization}): Formal-mathematische Optimierungsmodelle (0-1 ILP, Goal Programming, Credibility Programming) zur mathematischen Maximierung/Minimierung von Zielfunktionen.
        \item $v_{1,4}$ (\textit{Hybrid Soft Computing / Evidential Reasoning}): Mehrstufige Kombinationen aus Fuzzy Logic, Neuro-Fuzzy, DEA und Evidenztheorien (Dempster-Shafer).
    \end{itemize}

    \item \textbf{Parameter $P_2$: Mismatch Handling Phase (Integrations- und Mismatch-Handhabung)} \\
    Bestimmt den Zeitpunkt und die methodische Tiefe, mit der Fehlanpassungen (Mismatches) zwischen COTS-Funktion und Organisationsanforderung behandelt werden.
    \begin{itemize}
        \item $v_{2,1}$ (\textit{Pre-Evaluation / Requirements Phase}): Mismatches werden vorab im Requirements Engineering durch elastische Anforderungsverhandlung eliminiert.
        \item $v_{2,2}$ (\textit{Integrated Mismatch Resolution}): Mismatches und deren Behebungsaufwände (Adapter, Wrapping) fließen als explizite mathematische Variable direkt in die Evaluierungsrechnung ein.
        \item $v_{2,3}$ (\textit{Post-Evaluation Mismatch Management}): Die Auswahl erfolgt primär nach Funktionalität; Anpassungsstrategien werden erst nach Zuschlagserteilung konzipiert.
        \item $v_{2,4}$ (\textit{No Explicit Mismatch Handling}): Es findet kein explizites Mismatch-Handling statt; Produkte werden als starr angenommen.
    \end{itemize}

    \item \textbf{Parameter $P_3$: Criteria Catalog Standardization (Standardisierungsgrad des Kriterienkatalogs)} \\
    Beschreibt den Grad der Wiederverwendbarkeit und Strukturierung der zugrunde liegenden Bewertungskriterien.
    \begin{itemize}
        \item $v_{3,1}$ (\textit{Standardized International Frameworks}): Nutzung internationaler Taxonomien (z.\,B. ISO/IEC 25010, ISO 9126, IEEE 1061).
        \item $v_{3,2}$ (\textit{Domain-Specific Adaptations}): Auf spezifische Softwaregattungen (z.\,B. ERP, CRM, MES, IoT-Plattformen) zugeschnittene Standardkataloge.
        \item $v_{3,3}$ (\textit{Ad-hoc / Custom Criteria}): Rein projektspezifisch und frei formulierte Kriterienlisten ohne standardisiertes Metamodell.
    \end{itemize}

    \item \textbf{Parameter $P_4$: Uncertainty Handling Mechanism (Modellierung von Ungewissheit und Unschärfe)} \\
    Gibt an, wie Informationsunsicherheiten und vage Expertenurteile mathematisch verarbeitet werden.
    \begin{itemize}
        \item $v_{4,1}$ (\textit{Deterministic / Crisp Scoring}): Verwendung exakter Punktwerte (z.\,B. Skala 1--5) ohne Unschärfeberücksichtigung.
        \item $v_{4,2}$ (\textit{Fuzzy Logic \& Fuzzy Sets}): Transformation in unscharfe Zahlen (z.\,B. dreieckige/trapezförmige Fuzzy-Zahlen, Intuitionistic Fuzzy).
        \item $v_{4,3}$ (\textit{Stochastics \& Credibility Theory / Evidential Reasoning}): Stochastische Verteilungen, Chance-Constrained Credibility Verteilungen oder Dempster-Shafer Evidenzintervalle.
    \end{itemize}

    \item \textbf{Parameter $P_5$: Decision Support Automation Level (Automatisierungsgrad der Entscheidungsunterstützung)} \\
    Definiert das Ausmaß der werkzeuggestützten Software-Automatisierung des Auswahlprozesses.
    \begin{itemize}
        \item $v_{5,1}$ (\textit{Manual / Spreadsheet Calculations}): Manuelle Berechnung, z.\,B. via MS Excel oder Papier-Protokollen.
        \item $v_{5,2}$ (\textit{Dedicated Prototypic DSS Tooling}): Eigenständige Berechnungssoftware oder Prototypen zur Algorithmenausführung.
        \item $v_{5,3}$ (\textit{Knowledge-Based \& Expert Systems}): Vollautomatisierte Wissensdatenbanken, Regel-Engines oder intelligente Agenten-Systeme.
    \end{itemize}

    \item \textbf{Parameter $P_6$: Covered Criteria Scope (Abdeckungsbereich der Evaluierungskriterien)} \\
    Legt fest, welche Dimensionen der Softwarequalität und der Kontextbedingungen evaluiert werden.
    \begin{itemize}
        \item $v_{6,1}$ (\textit{Technical \& Functional Quality Only}): Ausschließliche Bewertung technischer und funktionaler Merkmale.
        \item $v_{6,2}$ (\textit{Technical \& Vendor/Business Quality}): Erweiterung um geschäftliche Kriterien (Herstellerstabilität, Lizenzkosten, Support).
        \item $v_{6,3}$ (\textit{Comprehensive Architectural \& Mismatch Metrics}): Ganzheitliche Abdeckung inklusive architektonischer Kopplung, Mismatch-Kosten und Integrationsaufwänden.
    \end{itemize}
\end{enumerate}

\subsubsection{Empirische Häufigkeitsverteilung der Parameterzustände}

\autoref{tab:morphological_frequencies} zeigt die absolute und relative Verteilung aller Wertzustände $v_{i,j}$ in der untersuchten Stichprobe von $N_{\text{reports}} = 41$.

\begin{table}[htbp]
\centering
\caption{Empirische Häufigkeitsverteilung der Parameterausprägungen in den 41 Primärstudien}
\label{tab:morphological_frequencies}
\begin{adjustbox}{max width=\textwidth}
\begin{tabular}{llrr}
\toprule
\textbf{Parameter} & \textbf{Diskrete Wertzustände ($v_{i,j}$)} & \textbf{Absolut ($n$)} & \textbf{Relativ (\%)} \\
\midrule
\multirow{4}{*}{$P_1$: Math. Decision Model Family} 
 & $v_{1,1}$: Single-Criterion / Heuristic Matching & 4 & 9,8\% \\
 & $v_{1,2}$: Multi-Criteria Decision Analysis (MCDA) & 14 & 34,1\% \\
 & $v_{1,3}$: Mathematical Programming / Optimization & 9 & 22,0\% \\
 & $v_{1,4}$: Hybrid Soft Computing / Evidential Reasoning & 14 & 34,1\% \\
\midrule
\multirow{4}{*}{$P_2$: Mismatch Handling Phase} 
 & $v_{2,1}$: Pre-Evaluation / Requirements Phase & 6 & 14,6\% \\
 & $v_{2,2}$: Integrated Mismatch Resolution & 11 & 26,8\% \\
 & $v_{2,3}$: Post-Evaluation Mismatch Management & 8 & 19,5\% \\
 & $v_{2,4}$: No Explicit Mismatch Handling & 16 & 39,0\% \\
\midrule
\multirow{3}{*}{$P_3$: Criteria Catalog Standardization} 
 & $v_{3,1}$: Standardized International Frameworks (ISO 25010) & 15 & 36,6\% \\
 & $v_{3,2}$: Domain-Specific Adaptations & 11 & 26,8\% \\
 & $v_{3,3}$: Ad-hoc / Custom Criteria & 15 & 36,6\% \\
\midrule
\multirow{3}{*}{$P_4$: Uncertainty Handling Mechanism} 
 & $v_{4,1}$: Deterministic / Crisp Scoring & 16 & 39,0\% \\
 & $v_{4,2}$: Fuzzy Logic \& Fuzzy Sets & 18 & 43,9\% \\
 & $v_{4,3}$: Stochastics, Credibility \& Evidential Reasoning & 7 & 17,1\% \\
\midrule
\multirow{3}{*}{$P_5$: Decision Support Automation Level} 
 & $v_{5,1}$: Manual / Spreadsheet Calculations & 17 & 41,5\% \\
 & $v_{5,2}$: Dedicated Prototypic DSS Tooling & 15 & 36,6\% \\
 & $v_{5,3}$: Knowledge-Based \& Expert Systems & 9 & 21,9\% \\
\midrule
\multirow{3}{*}{$P_6$: Covered Criteria Scope} 
 & $v_{6,1}$: Technical \& Functional Quality Only & 7 & 17,1\% \\
 & $v_{6,2}$: Technical \& Vendor/Business Quality & 19 & 46,3\% \\
 & $v_{6,3}$: Comprehensive Architectural \& Mismatch Metrics & 15 & 36,6\% \\
\bottomrule
\end{tabular}
\end{adjustbox}
\end{table}

---

\subsection{Ergebnisse der Cross-Consistency Assessment (CCA)}
\label{subsec:cca_ergebnisse}

Der theoretische Problemraum umfasst $N = 4 \times 4 \times 3 \times 3 \times 3 \times 3 = 1.296$ theoretisch denkbare Profil-Kombinationen. Zur Reduktion dieses Feldes auf konsistente, valide Konfigurationen wurde eine **Cross-Consistency Assessment (CCA)** nach Ritchey durchgeführt. Hierbei werden sämtliche Parameterpaare auf ihre mathematische, logische und empirische Vereinbarkeit hin geprüft.

Die Paarevaluierung teilt das morphologische Feld in drei Beziehungszustände:
\begin{itemize}
    \item \textbf{Valid ($V$):} Logisch konsistente und empirisch durch Studien belegte Konfigurationen.
    \item \textbf{Empirisch unbesetzt ($T$):} Theoretisch denkbare und mathematisch widerspruchsfreie Kombinationen, die jedoch in der aktuellen Literatur bisher nicht untersucht wurden (Forschungslücken).
    \item \textbf{Logisch unmöglich ($X$):} Inhärent inkonsistente Kombinationen, die aufgrund ontologischer oder mathematischer Widersprüche ausgeschlossen werden müssen.
\end{itemize}

\subsubsection{Konstitutive Domänen- und Ausschlussregeln}

Die systematische Reduktion des Gesamtraums basiert auf vier zentralen Ausschlussregeln ($R_1$ bis $R_4$):

\begin{description}
    \item[\textbf{Regel 1 ($R_1$): Optimierung vs. manuelle Ausführung}] \\
    \textit{Kombinatorische mathematische Optimierungsmodelle ($P_1 = v_{1,3}$) schließen eine rein manuelle Tabellenkalkulations-Auswertung ($P_5 = v_{5,1}$) strikt aus ($X$).} \\
    \textbf{Begründung:} Die Lösung von NP-schweren 0-1 ILP- oder Chance-Constrained Credibility-Problemen erfordert spezialisierte mathematische Solver (z.\,B. CPLEX, Gurobi) oder dedizierte Algorithmen. Eine Handausführung ist bei realen Problemgrößen praktisch unmöglich.

    \item[\textbf{Regel 2 ($R_2$): Crisp vs. Unscharf}] \\
    \textit{Hochgradig mathematische Unschärfe-Mechanismen ($P_4 \in \{v_{4,2}, v_{4,3}\}$) sind mit reinen Ad-hoc-Heuristiken ($P_1 = v_{1,1}$) logisch inkompatibel ($X$).} \\
    \textbf{Begründung:} Wenn Entscheider komplexe Fuzzy-Zahlen oder Credibility-Verteilungen zur Modellierung von Vagheit nutzen, erfordert dies mathematische Transformations- und Defuzzifizierungs-Operatoren. Der Einsatz einfacher Checklisten widerhebt den Zweck fortgeschrittener Unschärfe-Theorien.

    \item[\textbf{Regel 3 ($R_3$): Post-Evaluation Mismatch vs. reine Qualität}] \\
    \textit{Verfahren, die ein integriertes Mismatch-Handling beanspruchen ($P_2 \in \{v_{2,1}, v_{2,2}\}$), müssen zwingend architektonische Metriken abdecken ($P_6 = v_{6,3}$). Eine Beschränkung auf reine funktionale Standardkriterien ($P_6 = v_{6,1}$) ist inkonsistent ($X$).} \\
    \textbf{Begründung:} Mismatches treten vor allem an den Schnittstellen, Schnittstellen-Kopplungen und Komponenten-Architekturen auf. Werden architektonische Kriterien ausgeblendet, können Mismatch-Auflösungsaufwände mathematisch nicht bewertet werden.

    \item[\textbf{Regel 4 ($R_4$): DSS-Tool-Automatisierung vs. flache Kriterienlisten}] \\
    \textit{Vollautomatisierte, wissensbasierte Expertensysteme ($P_5 = v_{5,3}$) setzen semantisch strukturierte Kriterienkataloge ($P_3 \in \{v_{3,1}, v_{3,2}\}$) voraus. Unstrukturierte Ad-hoc-Kriterien ($P_3 = v_{3,3}$) sind mit automatisierter Regelverarbeitung unvereinbar ($X$).} \\
    \textbf{Begründung:} Ein Inferenz-Engine benötigt formal definierte Entitäten und Attributrelationen, um automatisiertes Reasoning oder Scraped-Feature-Matching durchzuführen.
\end{description}

Durch Anwendung der Regeln $R_1$--$R_4$ wird der theoretische Raum von 1.296 Kombinationen auf 184 mathematisch valide Konfigurationen reduziert. Davon entfallen genau 41 empirische Realisierungen auf die gesichtete Primärliteratur, während 143 valide Zellen als empirisch unbesetzte weiße Flecken identifiziert wurden.

---

\subsection{Typologisierung und vergleichende Analyse der 6 methodischen Archetypen}
\label{subsec:typologisierung_archetypen}

Aus der morphologischen Synthese und der Zuordnung der 41 Primärstudien lassen sich **sechs konstitutive methodische Archetypen (A bis F)** aggregieren. Diese Archetypen repräsentieren die vorherrschenden Lösungsmuster der wissenschaftlichen Praxis.

\subsubsection{Archetyp A: Klassische MCDM- \& Outranking-Frameworks}
\begin{itemize}
    \item \textbf{Morphologisches Profil-Muster:} $(v_{1,2}, v_{2,4}, v_{3,1}, v_{4,1}, v_{5,1}, v_{6,2})$
    \item \textbf{Mittlerer QA-Score:} 7,2 / 8,0 (Höchste Transparenz und methodische Rigorosität)
    \item \textbf{Repräsentative Studien:} Blin \& Tsoukiàs \cite{M3B2D3GD}, Shyur \cite{U5JPV4SF}, Kontio et al. \cite{LZQJ6M4X}, Chatzipetrou et al. \cite{UK2X9C4W}, Taherdoost \& Mohebi \cite{LPC92LAP}.
    \item \textbf{Methodische Merkmale:} Nutzt mathematische MCDA-Klassiker wie AHP, ANP, ELECTRE oder TOPSIS zur Kriteriengewichtung und Alternativen-Rangfolge. Kriterien basieren meist auf ISO/IEC-Standards.
    \item \textbf{Stärken:} Höchste methodische Nachvollziehbarkeit, mathematisch bewiesene Konsistencyprüfungen (z.\,B. Consistency Ratio bei AHP), universelle Anwendbarkeit.
    \item \textbf{Praktische Grenzen:} Hoher manueller Paarvergleichsaufwand bei großen Kriterienkatalogen ($O(n^2)$ Komplexität), vernachlässigt Mismatch-Resolutionskosten und dynamische Unsicherheiten.
\end{itemize}

\subsubsection{Archetyp B: Kombinatorische Multi-Objective Optimierung \& Ressourcen-Allokation}
\begin{itemize}
    \item \textbf{Morphologisches Profil-Muster:} $(v_{1,3}, v_{2,2}, v_{3,3}, v_{4,1}, v_{5,2}, v_{6,3})$
    \item \textbf{Mittlerer QA-Score:} 6,3 / 8,0
    \item \textbf{Repräsentative Studien:} Neubauer \& Stummer (Odin) \cite{ZQMG69VT}, Cortellessa et al. (0-1 ILP) \cite{VUVLUN85}, Cortellessa et al. (Atana) \cite{MZ8JHDYA}, Aggarwal et al. \cite{W8P543Y2}.
    \item \textbf{Methodische Merkmale:} Formuliert das Auswahlproblem als exaktes mathematisches Optimierungsproblem (0-1 ILP, Goal Programming), um Multikomponentensysteme bezüglich Zuverlässigkeit, Hardwarekosten und Integrationsaufwand zu optimieren.
    \item \textbf{Stärken:} Berechnet echte Pareto-optimale Systemkonfigurationen, beherrscht kombinatorische Explosionen bei der Multi-Komponenten-Auswahl.
    \item \textbf{Praktische Grenzen:} Erfordert exakte mathematische Parameter, die in der Praxis schwer schätzbar sind; hohes Abstraktionsniveau erschwert die Akzeptanz bei Entscheidern.
\end{itemize}

\subsubsection{Archetyp C: Zielorientierte \& Mismatch-Resolution Architektur-Modelle}
\begin{itemize}
    \item \textbf{Morphologisches Profil-Muster:} $(v_{1,1}, v_{2,1}, v_{3,2}, v_{4,1}, v_{5,1}, v_{6,3})$
    \item \textbf{Mittlerer QA-Score:} 6,8 / 8,0
    \item \textbf{Repräsentative Studien:} Maiden et al. (CARE) \cite{TLB687YL, 6RNULWVY}, Alves et al. (AGORA) \cite{8XNYS4M9}, Mohamed et al. (MiHOS) \cite{MC9UMXZZ, 9F79VAAB}, Grau et al. \cite{DAU8NSK8}.
    \item \textbf{Methodische Merkmale:} Verknüpft Requirements Engineering (Goal Graphs, i*-Modelle) direkt mit der COTS-Evaluierung. Mismatches werden als primäre Steuerungseinheit verstanden und durch Nachverhandlung aufgelöst.
    \item \textbf{Stärken:} Hohes Praxisverständnis für architektonische Schnittstellen und organisatorische Anpassungsbedarfe.
    \item \textbf{Praktische Grenzen:} Geringe mathematische Formalisierung der Auswertung; stark abhängig von der Verhandlungskompetenz der Stakeholder.
\end{itemize}

\subsubsection{Archetyp D: Fuzzy Mathematische Programmierung \& Credibility Modelle}
\begin{itemize}
    \item \textbf{Morphologisches Profil-Muster:} $(v_{1,3}, v_{2,4}, v_{3,1}, v_{4,3}, v_{5,2}, v_{6,2})$
    \item \textbf{Mittlerer QA-Score:} 6,7 / 8,0
    \item \textbf{Repräsentative Studien:} Mehlawat \& Gupta \cite{6YJHIVHG, B7ZTQCZZ}, Gupta et al. \cite{VRTXHQI7}, Mehlawat et al. \cite{JX4RDRDU}, Sadiq et al. \cite{P9SATBH5}, Kalantari et al. \cite{UFCEV52H}.
    \item \textbf{Methodische Merkmale:} Verknüpft mathematische Optimierung mit unscharfer Logik und Credibility-Theorie. Modelliert Kosten und Performanz als Fuzzy-Chance-Restriktionen.
    \item \textbf{Stärken:} Extrem robuste Handhabung von ungenauen, vagen oder probabilistischen Messdaten bei hoher mathematischer Strenge.
    \item \textbf{Praktische Grenzen:} Mathematisch hochgradig komplex; erfordert Spezialsoftware und mathematisches Expertenwissen.
\end{itemize}

\subsubsection{Archetyp E: Advanced Hybrid Soft Computing \& Evidential Reasoning}
\begin{itemize}
    \item \textbf{Morphologisches Profil-Muster:} $(v_{1,4}, v_{2,2}, v_{3,1}, v_{4,2}, v_{5,2}, v_{6,2})$
    \item \textbf{Mittlerer QA-Score:} 7,1 / 8,0 (Höchste methodische Güte im modernen Feld)
    \item \textbf{Repräsentative Studien:} Kaur et al. \cite{NZJEEPAC}, Sadiq et al. \cite{QNWK8SAP, JT5X3FIG, DFASRUVC, LNQHES84}, Alidrisi et al. (FMDBA) \cite{FVINCQ89}, Wani et al. \cite{8KF4HAZS}, Raj et al. (DEA-Fuzzy) \cite{YVWGHAWQ}, Saremi et al. \cite{LYMHTPXS}, Al-Qiaisi et al. \cite{5BBYQCJB}.
    \item \textbf{Methodische Merkmale:} Mehrstufige Methodenkaskaden (z.\,B. Fuzzy-AHP + PROMETHEE/VIKOR, DEA, Dempster-Shafer Evidential Reasoning). Kombiniert die Vorteile verschiedener Soft-Computing-Ansätze.
    \item \textbf{Stärken:} Kompensiert methodische Nachteile einzelner Verfahren; exzellente Eignung für komplexe, datenintensive Auswahlentscheidungen.
    \item \textbf{Praktische Grenzen:} Hohe Komplexität der Datenpipeline; aufwendige Implementierung dedizierter DSS-Tools erforderlich.
\end{itemize}

\subsubsection{Archetyp F: Wissensbasierte \& Evidenzbasierte DSS-Plattformen}
\begin{itemize}
    \item \textbf{Morphologisches Profil-Muster:} $(v_{1,4}, v_{2,3}, v_{3,2}, v_{4,2}, v_{5,3}, v_{6,3})$
    \item \textbf{Mittlerer QA-Score:} 6,7 / 8,0
    \item \textbf{Repräsentative Studien:} Wanyama \& Far \cite{CZZ6XQZJ, WBGP3ETU}, Morera et al. (DesCOTS) \cite{FCYW9EWJ}, Jadhav \& Sonar (Plato) \cite{CI86GRRG, 8VHUMEKY}, Sen et al. (HKBS) \cite{U4E7VXQU}, Upadhyay et al. (COTS-ESF) \cite{WZM7DUBY}, Lempert (IoT PRISM) \cite{RRU375GB}.
    \item \textbf{Methodische Merkmale:} Nutzung vollautomatisierter Softwareplattformen, Wissensdatenbanken, Multi-Agenten-Systeme oder Inferenz-Engines zur Speicherung von COTS-Attributen und automatischen Regelauswertung.
    \item \textbf{Stärken:} Höchster Automatisierungsgrad; ermöglicht nachhaltige Wiederverwendung von Evaluierungswissen über mehrere Projekte hinweg.
    \item \textbf{Praktische Grenzen:} Hohe Initialinvestition in den Aufbau und die Pflege der Wissensbasis; Gefahr veralteter Produktdatenbanken.
\end{itemize}

---

\subsubsection{Vollständige Klassifikationsmatrix aller 41 primären Literaturprofile}

\autoref{tab:morphological_profiles_full} stellt das vollständige morphologische Profil-Tupel, die Archetypen-Zuordnung sowie den bewerteten Quality Appraisal Score (QA-Score auf einer Skala von 0,0 bis 8,0) für alle 41 analysierten Studienberichte dar.

\begin{longtable}{p{2.0cm} p{2.6cm} p{2.8cm} p{2.7cm} p{2.2cm} p{1.0cm} p{2.5cm}}
\caption{Vollständige morphologische Profil-Klassifikation und Qualitätsbewertung aller 41 Primärstudien} \label{tab:morphological_profiles_full} \\
\toprule
\textbf{Zitation} & \textbf{Autor \& Jahr} & \textbf{Methodenname / Fokussierung} & \textbf{Morphologisches Profil-Tupel} & \textbf{Archetypen-Cluster} & \textbf{QA-Score} & \textbf{Evidenzzitat (Quelle)} \\
\midrule
\end{firsthead}
\toprule
\textbf{Zitation} & \textbf{Autor \& Jahr} & \textbf{Methodenname / Fokussierung} & \textbf{Morphologisches Profil-Tupel} & \textbf{Archetypen-Cluster} & \textbf{QA-Score} & \textbf{Evidenzzitat (Quelle)} \\
\midrule
\end{head}
\midrule
\multicolumn{7}{r}{\textit{Fortsetzung auf der nächsten Seite}} \\
\end{foot}
\bottomrule
\endlastfoot

\cite{M3B2D3GD} & Blin \& Tsoukiàs (2001) & ELECTRE Software Quality & $(v_{1,2}, v_{2,4}, v_{3,1}, v_{4,1}, v_{5,1}, v_{6,1})$ & Archetyp A & 7,0 & \enquote{Multi-criteria methodology contribution} \\
\cite{TLB687YL} & Maiden et al. (2002) & CARE Process & $(v_{1,1}, v_{2,1}, v_{3,2}, v_{4,1}, v_{5,1}, v_{6,3})$ & Archetyp C & 7,5 & \enquote{COTS acquisition with requirements} \\
\cite{8XNYS4M9} & Alves et al. (2003) & AGORA Goal Matching & $(v_{1,1}, v_{2,1}, v_{3,2}, v_{4,1}, v_{5,1}, v_{6,3})$ & Archetyp C & 7,0 & \enquote{Integrating goal-oriented requirement} \\
\cite{QDUNCXFC} & Kunda \& Brooks (2004) & STACE Framework & $(v_{1,2}, v_{2,4}, v_{3,3}, v_{4,1}, v_{5,1}, v_{6,2})$ & Archetyp A & 6,5 & \enquote{Socio-technical approach for COTS} \\
\cite{CZZ6XQZJ} & Wanyama \& Far (2005) & Multi-Agent DSS Frame & $(v_{1,4}, v_{2,3}, v_{3,2}, v_{4,2}, v_{5,3}, v_{6,3})$ & Archetyp F & 6,5 & \enquote{Towards providing decision support} \\
\cite{LZQJ6M4X} & Kontio et al. (2005) & ELECTRE Screening & $(v_{1,2}, v_{2,4}, v_{3,1}, v_{4,1}, v_{5,1}, v_{6,2})$ & Archetyp A & 7,0 & \enquote{Evaluating COTS components} \\
\cite{U5JPV4SF} & Shyur (2006) & ANP-TOPSIS Hybrid & $(v_{1,2}, v_{2,4}, v_{3,1}, v_{4,1}, v_{5,2}, v_{6,2})$ & Archetyp A & 7,5 & \enquote{COTS evaluation using modified TOPSIS} \\
\cite{6RNULWVY} & Maiden et al. (2006) & CARE Architecture & $(v_{1,1}, v_{2,1}, v_{3,2}, v_{4,1}, v_{5,1}, v_{6,3})$ & Archetyp C & 7,0 & \enquote{Requirements engineering for COTS} \\
\cite{MC9UMXZZ} & Mohamed et al. (2007a) & MiHOS Mismatch Model & $(v_{1,1}, v_{2,2}, v_{3,3}, v_{4,1}, v_{5,2}, v_{6,3})$ & Archetyp C & 7,5 & \enquote{Decision support for selecting COTS} \\
\cite{9F79VAAB} & Mohamed (2007b) & MiHOS Architecture & $(v_{1,1}, v_{2,2}, v_{3,3}, v_{4,1}, v_{5,2}, v_{6,3})$ & Archetyp C & 7,5 & \enquote{Comprehensive mismatch handling} \\
\cite{ZQMG69VT} & Neubauer \& Stummer (2007) & Odin Interactive MOCO & $(v_{1,3}, v_{2,2}, v_{3,3}, v_{4,1}, v_{5,2}, v_{6,3})$ & Archetyp B & 7,0 & \enquote{Interactive decision support} \\
\cite{FCYW9EWJ} & Morera et al. (2007) & DesCOTS Infrastructure & $(v_{1,4}, v_{2,3}, v_{3,1}, v_{4,1}, v_{5,3}, v_{6,3})$ & Archetyp F & 7,0 & \enquote{DesCOTS: A decision support system} \\
\cite{DAU8NSK8} & Grau et al. (2008) & DesCOTS Goal Quality & $(v_{1,1}, v_{2,1}, v_{3,1}, v_{4,1}, v_{5,2}, v_{6,2})$ & Archetyp C & 7,0 & \enquote{Goal-driven COTS selection} \\
\cite{VUVLUN85} & Cortellessa et al. (2008) & COSE 0-1 ILP Model & $(v_{1,3}, v_{2,2}, v_{3,3}, v_{4,1}, v_{5,2}, v_{6,3})$ & Archetyp B & 6,5 & \enquote{A multiobjective optimization approach} \\
\cite{MZ8JHDYA} & Cortellessa et al. (2011) & Atana Mismatch Optimization & $(v_{1,3}, v_{2,2}, v_{3,3}, v_{4,1}, v_{5,2}, v_{6,3})$ & Archetyp B & 6,0 & \enquote{Atana: Optimization of COTS selection} \\
\cite{CI86GRRG} & Jadhav \& Sonar (2009) & Hybrid MCDM Framework & $(v_{1,4}, v_{2,4}, v_{3,2}, v_{4,1}, v_{5,3}, v_{6,2})$ & Archetyp F & 6,5 & \enquote{Evaluating and selecting software} \\
\cite{8VHUMEKY} & Jadhav \& Sonar (2011) & Plato DSS System & $(v_{1,4}, v_{2,4}, v_{3,2}, v_{4,2}, v_{5,3}, v_{6,2})$ & Archetyp F & 7,0 & \enquote{Plato: A decision support system} \\
\cite{WBGP3ETU} & Wanyama \& Far (2008) & Automated Trade-Off Engine & $(v_{1,4}, v_{2,2}, v_{3,2}, v_{4,2}, v_{5,3}, v_{6,3})$ & Archetyp F & 6,5 & \enquote{Automated trade-off analysis} \\
\cite{U4E7VXQU} & Sen et al. (2010) & HKBS Expert System & $(v_{1,4}, v_{2,4}, v_{3,2}, v_{4,2}, v_{5,3}, v_{6,2})$ & Archetyp F & 6,5 & \enquote{A hybrid knowledge-based system} \\
\cite{6YJHIVHG} & Mehlawat \& Gupta (2014) & Fuzzy Chance-Constrained & $(v_{1,3}, v_{2,4}, v_{3,1}, v_{4,2}, v_{5,2}, v_{6,2})$ & Archetyp D & 6,5 & \enquote{Multiobjective COTS selection} \\
\cite{B7ZTQCZZ} & Mehlawat \& Gupta (2015) & Credibility Programming & $(v_{1,3}, v_{2,4}, v_{3,1}, v_{4,3}, v_{5,2}, v_{6,2})$ & Archetyp D & 7,0 & \enquote{COTS products selection using fuzzy} \\
\cite{VRTXHQI7} & Gupta et al. (2013) & Credibility Multiobjective & $(v_{1,3}, v_{2,4}, v_{3,1}, v_{4,3}, v_{5,2}, v_{6,2})$ & Archetyp D & 6,5 & \enquote{Credibility-based multiobjective} \\
\cite{W8P543Y2} & Aggarwal et al. (2014) & Cost-Reliability Model & $(v_{1,3}, v_{2,4}, v_{3,3}, v_{4,1}, v_{5,2}, v_{6,2})$ & Archetyp B & 6,0 & \enquote{Multi-attribute COTS selection} \\
\cite{WZM7DUBY} & Upadhyay et al. (2015) & DM-PT / COTS-ESF & $(v_{1,4}, v_{2,3}, v_{3,2}, v_{4,2}, v_{5,3}, v_{6,2})$ & Archetyp F & 6,5 & \enquote{COTS-ESF: Expert selection framework} \\
\cite{LYMHTPXS} & Saremi et al. (2015) & FAHP-TOPSIS Enterprise & $(v_{1,4}, v_{2,4}, v_{3,1}, v_{4,2}, v_{5,1}, v_{6,2})$ & Archetyp E & 7,0 & \enquote{Enterprise application selection} \\
\cite{JX4RDRDU} & Mehlawat et al. (2016) & Credibility Optimization & $(v_{1,3}, v_{2,4}, v_{3,1}, v_{4,3}, v_{5,2}, v_{6,2})$ & Archetyp D & 7,0 & \enquote{Credibility-based optimization model} \\
\cite{P9SATBH5} & Sadiq et al. (2016) & Fuzzy Preference Model & $(v_{1,3}, v_{2,4}, v_{3,3}, v_{4,2}, v_{5,1}, v_{6,1})$ & Archetyp D & 6,5 & \enquote{Fuzzy preference modeling} \\
\cite{JT5X3FIG} & Sadiq et al. (2017) & Intuitionistic Fuzzy Model & $(v_{1,4}, v_{2,4}, v_{3,3}, v_{4,2}, v_{5,1}, v_{6,1})$ & Archetyp E & 6,5 & \enquote{Intuitionistic fuzzy sets in COTS} \\
\cite{DFASRUVC} & Sadiq et al. (2020) & Entropy-Fuzzy Selection & $(v_{1,4}, v_{2,4}, v_{3,1}, v_{4,2}, v_{5,1}, v_{6,2})$ & Archetyp E & 7,0 & \enquote{Entropy-based fuzzy MCDM model} \\
\cite{QNWK8SAP} & Sadiq et al. (2021) & AHP-VIKOR Hybrid & $(v_{1,4}, v_{2,4}, v_{3,1}, v_{4,2}, v_{5,2}, v_{6,2})$ & Archetyp E & 7,0 & \enquote{AHP-VIKOR hybrid framework} \\
\cite{NZJEEPAC} & Kaur et al. (2020) & Delphi-FAHP-PROMETHEE & $(v_{1,4}, v_{2,4}, v_{3,1}, v_{4,2}, v_{5,2}, v_{6,2})$ & Archetyp E & 7,5 & \enquote{Multi-stage hybrid MCDM approach} \\
\cite{UK2X9C4W} & Chatzipetrou et al. (2019) & Empirical Requirements & $(v_{1,2}, v_{2,1}, v_{3,3}, v_{4,1}, v_{5,1}, v_{6,2})$ & Archetyp A & 7,5 & \enquote{Empirical evaluation of software} \\
\cite{5BBYQCJB} & Al-Qiaisi et al. (2021) & Hybrid Fuzzy-MCDM System & $(v_{1,4}, v_{2,4}, v_{3,3}, v_{4,2}, v_{5,1}, v_{6,2})$ & Archetyp E & 7,0 & \enquote{A hybrid fuzzy decision-making} \\
\cite{RRU375GB} & Lempert (2021) & IoT PRISM Plattform & $(v_{1,4}, v_{2,3}, v_{3,2}, v_{4,1}, v_{5,3}, v_{6,3})$ & Archetyp F & 7,0 & \enquote{IoT-Software-Plattformen: Methode} \\
\cite{UFCEV52H} & Kalantari et al. (2021) & Fuzzy Intra-Coupling & $(v_{1,3}, v_{2,2}, v_{3,3}, v_{4,2}, v_{5,2}, v_{6,3})$ & Archetyp D & 6,5 & \enquote{Optimal components selection based} \\
\cite{FVINCQ89} & Alidrisi et al. (2022) & FMDBA Distance Approach & $(v_{1,4}, v_{2,4}, v_{3,1}, v_{4,2}, v_{5,2}, v_{6,2})$ & Archetyp E & 7,0 & \enquote{Fuzzy modified distance-based} \\
\cite{8KF4HAZS} & Wani et al. (2023) & Neuro-Fuzzy Selection & $(v_{1,4}, v_{2,4}, v_{3,3}, v_{4,2}, v_{5,2}, v_{6,1})$ & Archetyp E & 7,0 & \enquote{Neuro-fuzzy decision support system} \\
\cite{YVWGHAWQ} & Raj et al. (2024) & Integrated DEA-Fuzzy & $(v_{1,4}, v_{2,4}, v_{3,1}, v_{4,2}, v_{5,2}, v_{6,2})$ & Archetyp E & 7,5 & \enquote{Integrating DEA and Fuzzy MCDM} \\
\cite{LNQHES84} & Sadiq et al. (2024) & IDS-Tool Dempster-Shafer & $(v_{1,4}, v_{2,4}, v_{3,1}, v_{4,3}, v_{5,2}, v_{6,2})$ & Archetyp E & 7,5 & \enquote{IDS-Tool: Interactive Decision Support} \\
\cite{LPC92LAP} & Taherdoost \& Mohebi (2025) & Critical MCDM Review & $(v_{1,2}, v_{2,4}, v_{3,1}, v_{4,1}, v_{5,1}, v_{6,2})$ & Archetyp A & 7,0 & \enquote{A critical examination of multi-criteria} \\
\end{longtable}
